% ============================================================
% 2D TFI cross-scale section — drop-in text for sn-article.tex
%
% Suggested placement:
%   * Results subsection: after "Generalization to frustrated and
%     competing-interaction spin chains", before "Learning physically
%     aligned quantum structures".
%   * Methods paragraph: append to "Quantum spin-chain Hamiltonians"
%     (or its own Methods subsection).
%   * Figure: figures2d/fig_tfi2d_crossscale.pdf
% ============================================================

\subsection{Extension to two-dimensional lattices}

All results presented so far concern one-dimensional chains, a setting in
which quasi-exact classical methods such as the density-matrix
renormalization group remain available for cross-checking. As a first step
beyond this regime, we test whether the measurement structures learned by
Diffushadow transfer to two-dimensional systems, where tensor-network
methods face rapidly growing entanglement costs. We consider the
transverse-field Ising model on an $L_x \times L_y$ square lattice with
periodic boundary conditions in both directions, at fixed width $L_x=4$ and
increasing length $L_y$. Sites are ordered row by row, so that a vertical
bond always connects tokens at a fixed relative offset in the measurement
sequence, independent of $L_y$; the relative positional encoding can
therefore reuse, at larger $L_y$, exactly the geometric relations observed
during training. Importantly, no architectural modification of any kind is
required: the two-dimensional geometry enters only through the training
data.

Diffushadow is trained on classical-shadow measurements of the $4\times3$
(12-qubit) and $4\times4$ (16-qubit) tori and evaluated on the $4\times5$
(20-qubit) and $4\times6$ (24-qubit) tori, which are never seen during
training. Figure~\ref{fig:tfi2d} compares site-averaged spin correlations
estimated from Diffushadow-generated measurements with exact
diagonalization, for horizontal nearest-neighbor, vertical
nearest-neighbor, and horizontal next-nearest-neighbor displacements. At
the trained size the generated estimates track the exact results across
both phases (mean squared error $5.6\times10^{-4}$ averaged over the three
observables), and the accuracy degrades only gradually at the unseen sizes
($1.3\times10^{-3}$ at 20 qubits and $3.0\times10^{-3}$ at 24 qubits).
In particular, the vertical correlations---which probe bonds separated by
eight token positions in the measurement sequence---are reproduced with
accuracy comparable to the horizontal ones at all sizes, indicating that
the model has learned the two-dimensional lattice geometry rather than a
one-dimensional proximity rule. The residual cross-scale error concentrates
near the critical region ($g_c^{2\mathrm{D}}\approx0.753$ for this
parametrization in the two-dimensional limit), where correlation lengths
are largest and finite-size statistics are most size-sensitive, and the
distinction between the $r{=}1$ and $r{=}2$ horizontal correlations is
partially washed out deep in the paramagnetic regime at the largest unseen
size. These results constitute a proof of principle that the cross-scale
generation mechanism underlying Diffushadow is not restricted to
one-dimensional physics, and they mark the regime---two-dimensional
lattices of growing extent---in which learned measurement generation may
eventually complement tensor-network approaches whose cost grows
exponentially with lattice width.

% ---------------- closing paragraph: tube thermodynamic limit ----------
% (add after the paragraph above; figure = fig14_tube_thermolimit)

Finally, we execute the complete inference pipeline in this
two-dimensional setting. Applying the recursive scale-extrapolation
procedure, the model is retrained on its own generated measurements and
extended to $4\times8$, $4\times10$, and $4\times12$ tori (32--48 qubits),
with every recursion step validated against DMRG at the same size; the
deviations remain below $0.05$ outside the crossover region and saturate,
rather than compound, with recursion depth. Extrapolating the generated
correlations to $1/N \to 0$ yields thermodynamic-limit estimates for the
infinite $4\times\infty$ tube, which we compare against numerically exact
infinite-system iDMRG values. Outside the crossover window the
extrapolations agree with the exact answers to within $0.008$--$0.06$
(within one to a few standard errors of the finite-size fits), while
inside the crossover the deviations grow to $0.27$, delimiting the
validity window of the procedure. To our knowledge this constitutes the
first end-to-end verification of generative thermodynamic-limit inference
against exact infinite-system results in a two-dimensional geometry ---
including a transparent map of where the inference degrades. The analysis
also shows that finite-size extrapolations of generated observables must
be reported with fit uncertainties: when correlations have already
converged in system size, an unweighted linear fit in $1/N$ amplifies
small per-size fluctuations into spurious intercept shifts.

% ---------------- figure ----------------
\begin{figure}[ht]
\centering
\includegraphics[width=0.98\textwidth]{fig/fig_tfi2d_crossscale.pdf}
\caption{Cross-scale measurement generation for the two-dimensional
transverse-field Ising model on $4\times L_y$ tori.
Site-averaged spin correlations $\langle Z_i Z_{i+\hat{x}}\rangle$,
$\langle Z_i Z_{i+\hat{y}}\rangle$, and
$\langle Z_i Z_{i+2\hat{x}}\rangle$ estimated from Diffushadow-generated
classical-shadow measurements (markers) are compared with exact
diagonalization (solid lines) as functions of the transverse-field
parameter $g$.
Left: $4\times4$ (16 qubits, included in training).
Middle: $4\times5$ (20 qubits, unseen).
Right: $4\times6$ (24 qubits, unseen).
The model is trained only on $4\times3$ and $4\times4$ measurement data.
Dotted vertical lines mark the two-dimensional critical point
$g_c^{2\mathrm{D}}\approx0.753$.}
\label{fig:tfi2d}
\end{figure}

% ---------------- Methods paragraph ----------------
% Append to Methods:

\subsection{Two-dimensional transverse-field Ising benchmark}
The two-dimensional benchmark uses the transverse-field Ising Hamiltonian
$H_{\mathrm{2D}}(g)=-(1-g)\sum_{\langle ij\rangle}\sigma^z_i\sigma^z_j
-g\sum_i \sigma^x_i$ on $L_x\times L_y$ square lattices with periodic
boundary conditions in both directions, where $\langle ij\rangle$ runs over
nearest-neighbor bonds. Sites are mapped to the one-dimensional token
sequence in row-major order, $s=rL_x+c$ for row $r$ and column $c$, so that
horizontal bonds correspond predominantly to adjacent site indices and
vertical bonds to a constant index offset $L_x$. Training data consist of
classical-shadow snapshots of the exact ground states of the $4\times3$ and
$4\times4$ tori, obtained by matrix-free Lanczos diagonalization, at 13
values of $g$ uniformly spaced in $[0.05,0.95]$ with $3000$ snapshots per
value and system size. The network architecture, tokenization, training
objective, and hyperparameters are identical to the one-dimensional
experiments (4 layers, 8 heads, hidden dimension 128, rotary positional
embeddings with $\theta=10^4$, 100 epochs). For evaluation, $2\times10^4$
snapshots are generated per parameter value with entropy-guided sequential
decoding (one site unmasked per denoising step); site-averaged $ZZ$
correlations are estimated with the same inverse-channel factors and
median-of-means aggregation ($K=10$) as in the one-dimensional analysis,
and compared with exact diagonalization at 10 values of $g$, including the
near-critical point $g=0.75$.
